\section*{Chapter \ref{ch:m1:the-sell-side} --- The Sell Side}

\begin{solution}{exo:m1:the-sell-side:1}
$30 \times 10^6 \times 0.015 \times \sqrt{3/3} = \$450\,000$: three days of
linear selling carry the risk of one day of the full position.
\end{solution}

\begin{solution}{exo:m1:the-sell-side:2}
$T = 500\,000/(0.1 \times 5\,000\,000) = 1$ day. $I = 0.7 \times 0.015 \times
\sqrt{0.1} = 0.33\%$.
\end{solution}

\begin{solution}{exo:m1:the-sell-side:3}
Fee: $0.035 \times 800 = \$28$ million. The issuer sold 20 million shares at
\$40 that the market valued at \$46 that evening: \$120 million left on the
table, received by the investors who were allocated shares --- the banks'
clients, which is why allocation is itself a \omterm{def:m1:the-sell-side:flow}{franchise} asset.
\end{solution}

\begin{solution}{exo:m1:the-sell-side:4}
$T = 3/(0.1\times 6) = 5$ days. $I = 0.7 \times 0.025 \times \sqrt{0.5} =
1.24\%$; risk $= 0.025\sqrt{5/3} = 3.23\%$; $d = 1.24 + 1.645 \times 3.23 =
6.55\%$, that is \$1.31 a share.
\end{solution}

\begin{solution}{exo:m1:the-sell-side:5}
$d_z = \sigma\sqrt{Q/V}\,\bigl(\kappa + z/\sqrt{3\pi}\bigr)$, and the bracket
does not depend on $Q$. Nine times the size: three times the discount, 6\%.
\end{solution}

\begin{solution}{exo:m1:the-sell-side:6}
(a) Financing $4.251 + 7.195 = \$11.4$ billion, 28\% of the division;
intermediation $\$19.6$ billion, 47\%; fees $\$9.3$ billion, 23\% (the small
remainder is other revenue). Financing is 37\% of the FICC and equities
total. (b) Financing is carry; intermediation is spread, with some position
P\&L; investment banking fees are fees.
\end{solution}

\begin{solution}{exo:m1:the-sell-side:7}
(a) The residual volatility is $0.02\sqrt{0.6} = 1.55\%$; risk
$1.55\% \times \sqrt{2/3} = 1.26\%$; $d = 0.63 + 1.645 \times 1.26 =
2.71\%$. (b) The hedge lowers the break-even discount by 61 basis points for
a cost of 2: it is worth it thirty times over, which is why every block is
hedged within seconds of the trade.
\end{solution}

\begin{solution}{exo:m1:the-sell-side:8}
Inside the model: the square-root rule was stated as insensitive to duration
\emph{within reasonable limits}; at half of the volume the dealer is the
market, its selling is visible to everyone, and realised impact rises well
above $\kappa\sigma\sqrt{Q/V}$, so the first term is no longer constant in
$\pi$. Outside the model: volume is not a given. Other participants detect
the seller and step back or sell ahead of it, so the volume against which
50\% is measured shrinks, and the price path acquires a drift against the
dealer that the driftless random walk excludes.
\end{solution}

\begin{solution}{pb:m1:the-sell-side:1}
\textbf{1.} \$100 million; 50\% of daily volume.
\textbf{2.} $T = 4/(0.1\times 8) = 5$ days.
\textbf{3.} $0.7\times 0.018\times\sqrt{0.5} = 0.89\%$, \$891\,000.
\textbf{4.} $0.018\sqrt{5/3} = 2.32\%$, \$2.32 million.
\textbf{5.} $d = 0.89 + 1.645\times 2.32 = 4.71\%$; bid \$23.82.
\textbf{6.} $T = 2.5$ days; $d = 0.89 + 1.645 \times 1.64 = 3.59\%$.
\textbf{7.} Residual volatility $0.018/\sqrt2 = 1.27\%$; risk $1.16\%$;
$d = 0.89 + 1.645\times 1.16 = 2.80\%$.
\textbf{8.} The discount less the impact: $1.91\%$ of \$100 million,
\$1.91 million.
\textbf{9.} $0.89 + 1.16 = 2.05\%$; it loses money with probability
$1-\Phi(1) = 16\%$.
\textbf{10.} Expected cost $0.89 + 0.04 = 0.93\%$; standard deviation
$0.018\sqrt{2.5/3} = 1.64\%$ (the fund does not hedge).
\textbf{11.} $2.80 - 0.93 = 1.87\%$, \$1.87 million: the fund buys the
removal of a \$1.64 million standard deviation and of all execution effort.
\textbf{12.} Her benchmark is the decision price; against it an agency
execution can end 3\% or 4\% adrift on a bad day, for which she answers
personally, while the risk bid's shortfall is known, agreed and final at
11:41.
\textbf{13.} Add 0.4\% (hedged blocks lose exactly the move relative to the
index): this is the adverse-selection cost of the client's flow.
\textbf{14.} $150\times10^6 \times 0.12 \times 0.10 \times 3/252 = \$21\,429$.
\textbf{15.} No: capital costs about 1\% of the expected profit. Risk, not
capital, prices a three-day position; capital binds for positions held for
months.
\textbf{16.} Bidding at the impact cost, a discount of 0.89\%, gives up the
whole expected profit of \$1.9 million and still has zero expected loss;
below that the desk pays for the relationship out of the \$3 million, and
should do so only if losing the block means losing the client.
\textbf{17.} $-\$3$ million on the stock, $+\$1$ million on the futures:
$-\$2$ million, already 1.7 standard deviations of the hedged unwind.
\textbf{18.} That the position was taken to meet a client's demand and that
the desk's inventory and risk limits are designed not to exceed the
reasonably expected near-term demands of its clients: a documented client
enquiry, limits calibrated on past flow, and an unwind consistent with them.
\textbf{19.} $2.80 + 0.40 = 3.20\%$, that is \textbf{80 cents a share}.
\textbf{20.} The client may sell at the bid or walk away, and it sells more
readily when it believes the price is about to fall: the dealer has written
a put whose exercise is correlated with bad news.
\end{solution}

\begin{solution}{iq:m1:the-sell-side:1}
The \omterm{def:m1:the-sell-side:sales-trader}{sales-trader} works the client's order in the market as agent: the
client owns every fill and the execution risk, and the \omterm{def:m1:the-sell-side:sales-trader}{sales-trader} is
judged on execution quality against a benchmark. The trader buys or sells
against the client as principal: once the price is agreed the P\&L of the
position belongs to the trader's book.
\iqlookfor{principal versus agent stated in terms of who owns the risk after
the phone call ends.}
\end{solution}

\begin{solution}{iq:m1:the-sell-side:2}
A linearly shrinking position has the variance of a constant one held a
third as long: its standard deviation is $\sqrt{1/3} \approx 58\%$ of that of
holding everything for nine days --- the risk of holding it all for three.
\iqlookfor{the $T/3$ rule from memory or derived in a minute.}
\end{solution}

\begin{solution}{iq:m1:the-sell-side:3}
How big is it against daily volume and what is the volatility (the
mechanical discount); can I hedge it, and with what; why is this client
trading, and what happened after its previous blocks; is there news or an
event (earnings, index change) within my unwind window; what is my current
position and limit in the name and the sector; and what is the client worth
to the bank.
\iqlookfor{\omterm{def:m1:what-a-trading-firm-does:adverse-selection}{adverse selection} and existing inventory mentioned, not only the
formula.}
\end{solution}

\begin{solution}{iq:m1:the-sell-side:4}
Proprietary trading was prohibited to banks and market-risk capital
requirements rose, so slow directional risk earns too little per unit of
capital. Financing is client-driven and collateralised, consumes little
market-risk capital per dollar of revenue, recurs every day, and grew with
the hedge-fund industry that the banks' own former traders went to build.
\iqlookfor{return on capital as the organising idea, plus the regulatory
cause.}
\end{solution}

\begin{solution}{iq:m1:the-sell-side:5}
One block of 4 million costs $2c$ per share if 1 million costs $c$: a total
of $8c$ million. Four blocks of 1 million cost $4c$ million: half as much
--- \emph{if} impact is fully forgotten from one day to the next and the
market does not infer that three more blocks follow. If impact is permanent
or the programme is detected, the later blocks are priced off a lower price
and the advantage shrinks toward zero. The client also carries the price
risk of the unsold shares for three more days.
\iqlookfor{the concavity argument, then the two caveats (impact decay,
information leakage) without being asked.}
\end{solution}

\begin{solution}{iq:m1:the-sell-side:6}
$\Var\bigl(\tfrac1T\int_0^T W_t\dd t\bigr) = \tfrac{1}{T^2}\iint
\min(s,t)\dd s\dd t = \tfrac{2}{T^2}\int_0^T\!\int_0^t s\dd s\dd t =
\tfrac{2}{T^2}\cdot\tfrac{T^3}{6} = T/3$ (times $\sigma^2$).
\iqlookfor{the covariance $\min(s,t)$ and a clean double integral; bonus for
the integration-by-parts route, $\int_0^T (T-t)\dd W_t$.}
\end{solution}
